
\section{Mathématiques : mode d'emploi}
\subsection{But}

Notre objectif en mathématiques est d'établir de manière logique des propriétés sur des objets en partant de : 
\begin{itemize}
    \item propriétés admises (axiomes)
    \item propriétés déjà établies
\end{itemize}

\subsection{Support matériel}

Les objets mathématiques sont représentés par des lettres, parfois déclinées avec des 
"machins" bizarres ou indexées. 

\uline{Règle fondamentale :} deux objets différents portent deux noms différents. 

\subsection{Objets étudiés}

En mathématiques actuelles, on étudie les ensembles. 

\uline{Relation fondamentale entre ensembles :} l'appartenance. Si $A$ et $B$ sont des 
ensembles, deux cas possibles : 

\begin{itemize}
    \item $A$ appartient à $B$, noté $A \in B$
    \item $A$ n'appartient pas à $B$, noté $A \notin B$
\end{itemize}

\begin{definition}{Inclusion}{}
    Si $A$, $B$ sont des ensembles, on dit que $A$ est inclus dans $B$ et on note 
    $A \subset B$ ssi pour tout élément $x$ de $A$, on a $x \in B$.
\end{definition}

\subsection{La théorie des ensembles}

\underline{Problème :} il faut choisir une axiomatique assez riche pour être utile, 
pas trop pour ne pas être contradictoire. 

\underline{Choix :} Théorème des ensembles de Zermelo-Fraenkel avec l'axiome du choix.

C'est-à-dire, avec des mots simples : 

\begin{enumerate}
    \item Deux ensembles sont identiques ssi ils ont les même éléments. 
    \item Il existe un ensemble, l'ensemble vide, ne contenant aucun élément.
    \item Si $A$ et $B$ sont des ensembles, il existe un ensemble $C = \{A, B\}$
            dont les éléments sont exactement $A$ et $B$.
    \item Réunion : si $A$ est un ensemble, il existe un ensemble $B$ dont les éléments
        sont les éléments des éléments de $A$. Donc $B = \bigcup_{x \in A} x$.
    \item Ensemble des parties : si $A$ est un ensemble, il existe un ensemble noté
        $\mathbb P(A)$ dont les éléments sont les sous-ensembles (parties) de $A$.
    \item Produit cartésien : Si $A$ et $B$ sont deux ensembles, il existe un ensemble
        noté $A \times B$ dont les éléments sont les couples $(x, y)$ où $x \in A$ 
        et $y \in B$.
    \item Axiome de compréhension : Si $A$ est un ensemble et $P$ une propriété réalisable
        par des éléments de $A$, on peut constituer l'ensemble $B = \{x \in A \mid P(x)\}$.
    \item Axiome de l'infini : L'ensemble $\mathbb N$ existe.
    \item Axiome du choix : Si $A$ est un ensemble d'ensembles non vides, il existe une 
        application $\fonction{f}{A}{E}{P}{x}$ avec $x \in P$ où $E = \bigcup_{P \in A} P$.
\end{enumerate}

\section{Les types de raisonements}

\subsection{Règle générale}

C'est le réusltat à démontrer qui guide dans le choix du raisonnement utilisé et des
acteurs mis en oeuvre. 

\subsection{Types d'énoncés}
\begin{enumerate}
	\setulcolor{vorange}
    
	\item\ul{Cas quelconques :}
    raisonnement par l'absurde (à utiliser en dernier recours).

    \item\ul{Énoncés du type $P$ ou $Q$ :}
    on suppose que $P$ n'est pas réalisé et on montre $Q$.

    \item\ul{Énoncés $P$ implique $Q$ :}
    \begin{itemize}
        \item On suppose $P$, on en déduit par une suite de donc que $Q$ l'est aussi.
        \item On peut également raisonner par contraposée, supposer non Q et en déduire non P.
    \end{itemize}

    \item\ul{Énoncés $P$ équivalent à $Q$ :}
    on montre les deux implications (raisonnement par double implication).

    \item\ul{Énoncés $A = B$ :}
    on montre les deux inclusions (raisonnement par double inclusion).

    \item\ul{Énoncés du type pour tout $x$ de $E$, $P(x)$ :}
    on considère un élément $x$ de $E$ et on établit $P(x)$.

    \item\ul{Énoncés du type il existe un unique $x$ dans $E$, $P(x)$ :}

    \textbf{Analyse - synthèse} :
    \begin{itemize}
        \item Analyse : On suppose qu'un tel $x$ existe et on l'identifie. On en 
        déduit l'unicité en cas d'existence. 
        \item Synthèse : on vérifie que la valeur identifiée convient. 
    \end{itemize}


    \textbf{Existence - unicité} :
    \begin{itemize}
        \item Exitence : on exhibe un $x$ qui convient
        \item Unicité : on suppose que $x_1$ et $x_2$ conviennent, et on montre que 
        $x_1 = x_2$
    \end{itemize}

    \item\ul{Énoncés du type il existe $x$ dans $E$, $P(x)$ :}
    on devine et exhibe un $x$ qui convient ou on commence une analyse jusqu'à obtenir
    suffisament d'informations pour identifier un tel $x$.

    \item\ul{Enoncé pour tout $n$ dans $\N$, $P(n)$ :}
    on peut raisonner par récurrence : 
    \begin{itemize}
        \item simple
        \item double
        \item avec prédécesseurs (= forte)
    \end{itemize}
\end{enumerate}
